%! The Idea of Elimination — Unit 2, Lesson 2.1 — Lesson Plan
\documentclass[10pt]{article}
\usepackage{linalg-article}
\usepackage{linalg-boxes}
\usepackage{graphicx}   % -article does NOT load graphicx; needed for thumbnails/screenshots
\graphicspath{{images/}}
\newcommand{\TallMath}[1]{$\displaystyle #1\rule[-1.4em]{0pt}{3.2em}$}
\newcommand{\xx}{\mathbf{x}}
\newcommand{\bb}{\mathbf{b}}
\newcommand{\UnitNumberName}{Unit 2: Solving Linear Equations Ax = b \quad}
\newcommand{\LessonNumberName}{Lesson 2.1: The Idea of Elimination}

\begin{document}
\begin{center}
    {\Large\bfseries \CourseName: \SchoolYear \\
    {\small\bfseries \UnitNumberName \LessonNumberName}}
\end{center}

\begin{tcolorbox}[colback=blush, colframe=burgundy, boxrule=0.9pt, arc=2mm,
    left=2mm, right=2mm, top=1.5mm, bottom=1.5mm]
    \textbf{Primary Objective:} Students can solve a $2\times 2$ (and a small $3\times 3$) system
    by \emph{elimination}: pick a \emph{pivot}, compute the \emph{multiplier} $\ell = \text{(entry
    to clear)}\div\text{(pivot)}$, subtract $\ell$ times the pivot row to make a zero, reach
    \emph{upper-triangular} form, and finish by \emph{back-substitution} --- then check by
    rebuilding $\bb$. They can explain why a zero in a pivot position calls for a \emph{row
    exchange}, and connect a breakdown to no solution / infinitely many.
\end{tcolorbox}

\renewcommand{\arraystretch}{1.4}
\begin{skillbox}[Priority Ideas \& Skills]{goldbox}
\begin{tabularx}{\linewidth}{@{}X|X@{}}
\textbf{Priority Ideas \& Skills} & \textbf{Key Understandings (the ``why'')} \\[2pt]
\hline
\vspace{-6pt}
\begin{itemize}[leftmargin=1.1em, nosep, after=\vspace{2pt}]
  \item Identify the \textbf{pivot} and compute the \textbf{multiplier} $\ell=$ entry$\,\div\,$pivot.
  \item Subtract $\ell\times$(pivot row) to make a chosen entry zero.
  \item Reach \textbf{upper-triangular} form, then \textbf{back-substitute}.
  \item Handle a \textbf{zero pivot} with a row exchange; read a breakdown as none / infinitely many.
\end{itemize}
&
\vspace{-6pt}
\begin{itemize}[leftmargin=1.1em, nosep, after=\vspace{2pt}]
  \item In 2.0 we ``subtracted to kill a variable''; the multiplier makes that a \emph{reliable rule}.
  \item The multiplier says \emph{exactly} how many pivot rows to subtract to zero out an entry.
  \item Triangular form is elimination's finish line --- solvable from the bottom up.
  \item A zero pivot is a road sign, not a dead end: swap, or read off no/infinite solutions.
\end{itemize}
\end{tabularx}
\end{skillbox}

\begin{skillbox}[{Vocabulary, Concepts, \& Theorems}]{greenbox}
\renewcommand{\arraystretch}{1.3}
\begin{tabularx}{\linewidth}{@{}l X@{}}
\textbf{Pivot} & The first nonzero coefficient in a row; we use it to clear the entries below it. \\
\textbf{Multiplier $\ell$} & $\ell = \dfrac{\text{entry to eliminate}}{\text{pivot}}$; how many pivot rows to subtract. \\
\textbf{Elimination step} & Replace a lower row by (that row) $-\;\ell\times$(pivot row) to make a zero. \\
\textbf{Upper-triangular system $U$} & All zeros below the diagonal: each equation has one fewer unknown going down. \\
\textbf{Back-substitution} & Solve the bottom equation, then work upward one unknown at a time. \\
\textbf{Row exchange} & Swap two equations to move a nonzero into a zero pivot position (formalized in 2.4). \\
\textbf{Breakdown} & A missing pivot: fixed by a swap, or a sign of no solution ($0=$ nonzero) / infinitely many ($0=0$). \\
\end{tabularx}
\end{skillbox}

\begin{fixedskillbox}[Activate Prior Knowledge \& Spiral Review]{blush}
    The Warm-Up rehearses what elimination is built from:
    \textbf{(1)} solving a system in which the variables \emph{already} line up (a direct subtract,
    the Lesson~2.0 move), \textbf{(2)} scaling an equation by a constant (multiplying both sides),
    and \textbf{(3)} writing $A\xx=\bb$ as two row equations (Lesson~2.0). Debrief item~1 out loud:
    it works \emph{because} the coefficients match --- today's job is what to do when they
    \emph{don't}.
\end{fixedskillbox}

\begin{skillbox}[Hook]{blush}
    Put two systems side by side. Left: $x+2y=5,\ 3x+2y=11$ --- both have a $2y$, so one
    subtraction kills $y$ (that was Lesson~2.0). Right: $2x+3y=13,\ 4x+7y=27$ --- now the
    $x$-terms are $2$ and $4$, so a plain subtraction leaves \emph{both} variables. Ask:
    \textbf{``We can't just subtract --- what could we do to the first equation \emph{first} so a
    subtraction wipes out $x$?''} Draw out ``double it.'' Name the number you doubled by --- the
    \textbf{multiplier} --- and you have the whole idea of elimination.
\end{skillbox}

\begin{skillbox}[Lesson]{blush}
\begin{multicols}{2}
\textbf{1. Why 2.0's trick needs an upgrade.} Last lesson both equations happened to share a term,
so subtracting once removed a variable. Usually the coefficients \emph{don't} match. Elimination
is the systematic fix: scale first, then subtract.

\textbf{2. Pivot and multiplier.} Take $2x+3y=13,\ 4x+7y=27$. The \textbf{pivot} is the first
coefficient we build on: the $2$ in front of $x$ in row~1. To clear the $4$ below it, ask ``$4$ is
how many $2$s?'' --- that is the \textbf{multiplier} $\ell=4\div2=2$.

\textbf{3. The elimination step.} Replace row~2 by row~2 $-\,\ell\times$row~1:
$(4x+7y)-2(2x+3y)=27-2(13)$ gives $y=1$. The $x$-term is gone \emph{by design} --- that's what the
multiplier guarantees.

\columnbreak

\textbf{4. Triangular, then back up.} We now have the \textbf{upper-triangular} system
$2x+3y=13$, $\;y=1$. \textbf{Back-substitute}: put $y=1$ into the top row, $2x+3=13$, so $x=5$.
\textbf{Check} by rebuilding $\bb$: $5\begin{bsmallmatrix} 2\\4 \end{bsmallmatrix}+1\begin{bsmallmatrix} 3\\7 \end{bsmallmatrix}=\begin{bsmallmatrix} 13\\27 \end{bsmallmatrix}$.

\textbf{5. Bigger systems, same staircase.} For $3\times 3$, use row~1's pivot to clear the whole
first column, then row~2's new pivot to clear below it. Each step makes one more zero; you march
down to a triangle, then back-substitute up.

\textbf{6. When the pivot is $0$.} If a pivot slot holds a $0$ but a lower row has a nonzero there,
\textbf{exchange rows} and continue. If nothing can be swapped in, elimination stalls: $0=$ a
nonzero means \emph{no} solution; $0=0$ means \emph{infinitely many}. Full row-swapping is 2.4.
\end{multicols}
\end{skillbox}

\begin{skillbox}[Active Monitoring]{redbox}
Circulate during the notes practice and Tier work. Watch for and correct:
\begin{itemize}[leftmargin=1.2em, nosep]
  \item multiplier upside-down --- it is entry$\,\div\,$pivot, not pivot$\,\div\,$entry;
  \item subtracting the left sides but forgetting to subtract $\ell\times$ the right-hand side too;
  \item back-substituting into the \emph{wrong} row, or quitting before checking by rebuilding $\bb$;
  \item panicking at a $0$ pivot instead of scanning below for a row to swap up.
\end{itemize}
Cold-call prompts: ``What's the pivot here?'' ``$\ell$ is how many pivot rows?'' ``Which entry did
that step zero out?'' ``How would you fix a $0$ in the pivot spot?''
\end{skillbox}

\begin{skillbox}[Group Work \& Differentiation]{redbox}
\textbf{Elimination Assembly Line} activity (tiered; mirrors the notes).
\begin{multicols}{3}
\textbf{Tier R --- Remediate.} One $2\times 2$ system with the multiplier already named: compute
$\ell$, do the elimination step, reach triangular form, back-substitute, and check.

\columnbreak
\textbf{Tier A --- Approaching.} A snack-pack context: set up $A\xx=\bb$, pick the pivot and
multiplier yourself, solve, and \emph{interpret} the counts. Then a $3\times 3$ staircase.

\columnbreak
\textbf{Tier E --- Extension.} A system needing a \textbf{row exchange} (zero pivot), plus one that
\emph{breaks down} --- decide no vs.\ infinitely many and justify from the equations.
\end{multicols}
\end{skillbox}

\begin{skillbox}[Individual Work \& Assessment]{redbox}
\textbf{Exit Ticket} (independent, no notes): solve one $2\times 2$ system by elimination ---
naming the pivot and multiplier, showing the triangular form, back-substituting, and checking ---
plus a \textbf{conceptual/justification check}: a system with a $0$ in the pivot position, where
students must say what move fixes it and why. Collect and sort into ``got the procedure / multiplier
slip / stuck on the zero pivot'' to decide whether 2.2 opens with a quick re-teach.
\end{skillbox}

\begin{skillbox}[Reinforcement \& Extension]{goldbox}
\textbf{Homework:} compute multipliers, run single elimination steps, solve $2\times 2$ and one
$3\times 3$ system to triangular form with a check, one applied snack/blend problem with an
interpret part, and a zero-pivot ``fix-it'' item. \textbf{Extension:} track the multipliers of a
$3\times 3$ elimination and notice they record the whole process --- a first look at the $L$ in
$A=LU$. \textbf{Preview:} Lesson~2.2, \emph{Elimination Matrices and Inverse Matrices} --- each
elimination step is itself a matrix, and undoing them leads to the inverse.
\end{skillbox}
% ── Teacher notes (migrated from the component keys) ──
% One per component, titled for it. They live here rather than in the _key
% files so that every component is the same length blank and keyed.

\begin{teachernote}[Warm-Up]
Item~1 is exactly the Lesson~2.0 move --- it works only \emph{because} both equations carry a
$2y$. Item~2 is the new ingredient: scaling a row so coefficients line up. Item~3 is the very
system solved in the notes; point out in the debrief that its $x$-coefficients ($2$ and $4$) do
\emph{not} match, so item~1's plain subtraction won't work --- that is why we need a multiplier.
\end{teachernote}

\begin{teachernote}[Guided Notes]
The heart of the lesson is the multiplier: $\ell=$ entry $\div$ pivot, applied as (lower row)
$-\,\ell\times$(pivot row). Watch for the flipped ratio (pivot $\div$ entry) and for students
forgetting to subtract $\ell\times$ the right-hand side. \S3's $3\times3$ is a big time sink ---
consider doing Round~1 with the class and letting groups finish Round~2. \S4 (zero pivot) is only a
preview of the permutations in 2.4.
\end{teachernote}

\begin{teachernote}[Group Activity]
Tier~A(1) is the interpret moment: whole positive pack counts make the algebra a real order. If a
group divides pivot by entry in Tier~A(1), they will get $\ell=\tfrac13$ --- remind them $\ell=$
entry $\div$ pivot, and that choosing the row-1 pivot of $1$ keeps it whole. Tier~E connects the
zero-pivot swap (fixable) with a genuine breakdown ($0=$ nonzero vs.\ $0=0$), the same
no/infinitely-many split from Lesson~2.0.
\end{teachernote}

\begin{teachernote}[Exit Ticket]
Item~1 checks the full procedure; full credit needs the pivot/multiplier named \emph{and} a check.
Item~2 is the conceptual check --- students must name the \emph{row exchange} and say \emph{why}
(no nonzero pivot without it), not just solve. Sort responses into ``procedure solid / multiplier
slip / stuck on the zero pivot'' to set up Lesson~2.2.
\end{teachernote}

\end{document}
